% ! TeX root = distributed-systems-final-report.tex
\section{Self-evaluation}\label{self-evaluation}

This project has proven itself an engaging testing ground to analyze some of the
challenges of distributed systems in a simulated real-world, tangible scenario.
Many architectural choices have been taken (section \ref{design}) and techniques
have been studied and applied (behaviour section, \ref{arch:behaviour}) trying
to address project's requirements and common distributed systems problems, while
other possible methods have been only discussed for future studies.

Observing the simulation, cars behave almost as expected: the traffic flows
fluidly, especially with the use of the controller; cars yield, slow down or
accelerate when they should and maintain the proper safety distance; vehicles in
failsafe mode are more cautious; the controller can identify ghosts with
sufficient precision when there aren't too many disconnected vehicles, so that
it can be aware of them in its calculations. Even if a vehicle is abruptly set
to disconnected, this still won't cause accidents.
\\
\\
However, the implementation has revealed itself trickier than expected. There
are still bugs or imprecisions in managing the physics and the collisions.
Therefore there are more collisions than we'd have hoped for, especially with
many vehicles in failsafe mode, indicating that the failsafe logic may most of
all require more improvements.

Some example results (from the testing section, \ref{validation-acceptance}),
coming from averaging multiple 10 minutes simulation runs and roughly rounding
their results to get an idea, using 4 roads, are: about 1 collision per minute
with 40 vehicles in normal mode, 2 per minute with 20 vehicles, 0.3 with 10
vehicles; about 1 collision every 10 seconds with 40 vehicles, of which 10 to 30
are in failsafe or disconnected mode. Moreover, one run with 40 vehicles out of
40 in a disconnected state ended up in a deadlock inside the roundabout (same
scenario described in \ref{arch:behaviour_controller}).

The traffic flow as well doesn't seem to change much with or without the
controller's help, since the total number of spawned vehicles is always around
300, for the same 40 vehicles test configurations as above.
\\
\\
The trickiest part is always to manage the collisions between a vehicle about to
enter the roundabout and one inside. Of course, one could have set the entering
vehicle to always stop and only enter if the roundabout is empty or almost
empty, but that would have made the project not very interesting: the main point
was exactly to try to make the vehicles coordinate with each other, while
maintaining the safety and optimizing the overall traffic flow.

Given the available time for the development, correcting the collisions bugs
would have meant neglecting many of the distributed systems aspects, so this way
of managing the project has been chosen for its higher educational value and
stimulating appeal. Regarding the controller's logic, then, surely a more
efficient evaluation process would have made the traffic more efficient, but
algorithms are yet another subject, not worth spending time on, in the context
of a small distributed systems project: the implemented controller's checks
were the sole ones sufficient to make the traffic flow and, then, to address
the distributed systems problems.
\\
\\
Another weak point has of course been testing, widely discussed in
\ref{validation}.
